\section{Microservicios y modelo de datos}
\label{sec:microservicios}
\textit{Responsables: cada dueño escribe la subsección de su microservicio
(archivo en \texttt{secciones/microservicios/}).}

Los 21 servicios lógicos del diseño se agrupan en seis microservicios desplegables, uno por
repositorio \texttt{nubeduca-ms-*}. Todos parten de la plantilla de servicio (LIB-01): outbox con
validación de esquemas, inbox idempotente y \texttt{serverless.yml} base.

\begin{table}[htbp]
\centering\small
\begin{tabularx}{\textwidth}{p{0.20\textwidth}p{0.15\textwidth}Xp{0.16\textwidth}}
\toprule
Microservicio & \texttt{source} & Publica & Responsable \\
\midrule
ms-usuarios & \texttt{lms.users} & \texttt{tenant.created}, \texttt{term.started}, \texttt{user.provisioned} & Benjamín, Edinson \\
ms-cursos & \texttt{lms.courses} & \texttt{course.created}, \texttt{course.published}, \texttt{course.archived}, \texttt{enrollment.created}, \texttt{enrollment.dropped}, \texttt{announcement.published} & Benjamín, Edinson \\
ms-contenidos & \texttt{lms.content} & \texttt{file.scanned}, \texttt{module.published} & Alexander \\
ms-evaluaciones & \texttt{lms.assessment} & \texttt{assignment.published}, \texttt{submission.submitted}, \texttt{submission.graded} & Edinson \\
ms-calificaciones & \texttt{lms.grading} & \texttt{grade.changed}, \texttt{grade.posted} (bus + SNS FIFO) & Edinson, Benjamín, Alexander \\
ms-notificaciones & --- & (solo consume) & Guillermo \\
\bottomrule
\end{tabularx}
\caption{Los seis microservicios, su \texttt{source} en EventBridge y los eventos que publican.}
\label{tab:microservicios}
\end{table}

\subsection{Convenciones comunes de datos}
Cada microservicio usa una sola tabla DynamoDB llamada \texttt{<stage>-t\_<nombre>} (modo
\textit{on-demand}), con clave de partición \texttt{tenant\_id} y clave de ordenamiento \texttt{sk}.
Además de sus entidades, la plantilla agrega tres tipos de ítem técnicos:

\begin{table}[htbp]
\centering\small
\begin{tabularx}{\textwidth}{lX}
\toprule
\texttt{sk} & Uso \\
\midrule
\texttt{EVENTO\#<eventId>} & Outbox: \texttt{source}, \texttt{detailType}, \texttt{detail} (JSON) y \texttt{ttl} de 7 días. Lo lee el Stream. \\
\texttt{INBOX\#<eventId>} & Inbox del consumidor: marca el evento como procesado (idempotencia). \\
\texttt{AGG\#<tipo>\#<id>} & Última \texttt{aggregateVersion} vista, para descartar eventos viejos. \\
\bottomrule
\end{tabularx}
\caption{Ítems técnicos de la plantilla en todas las tablas.}
\end{table}

Cada \texttt{poblar\_datos.py} carga 30\,000 registros en la tabla principal: UTEC 15\,000, UPC
10\,000 y PUCP 5\,000, sin escribir ítems \texttt{EVENTO\#} para no inundar el bus.

\subsection{ms-usuarios --- Tenants, periodos, usuarios y perfiles}
\label{sec:ms-usuarios}
\textit{Responsable: Benjamín (ACA-01) y Edinson (ACA-02). Repositorio: \url{https://github.com/nubeduca/nubeduca-ms-usuarios}.}

ms-usuarios es la base del resto: guarda las instituciones (tenants), sus periodos académicos y el perfil de cada usuario. Los tres tenants (UTEC, UPC y PUCP) los crea el seed y emite un \texttt{tenant.created} por cada uno. El administrador crea periodos en \texttt{PLANNED} y los inicia (\texttt{ACTIVE}); al iniciarlos se emite \texttt{term.started}, y ms-cursos solo acepta cursos en periodos activos. El nombre de un periodo es único en su tenant gracias a un ítem de reserva \texttt{TERMNAME\#} que se escribe en la misma transacción.

El perfil se crea en el primer login: el frontend llama a \texttt{POST /users/me} justo después de autenticarse. El tenant, el rol y el \texttt{userId} salen del token de Cognito, nunca del cuerpo; el correo y el nombre salen de los \textit{claims}, y lo que falte, junto con el código institucional, sale del padrón \texttt{ROSTER\#<email>} que la institución carga de antemano. La llamada es idempotente: la primera vez responde 201 y emite \texttt{user.provisioned}; las siguientes devuelven el perfil existente con 200 y sin evento. El evento no lleva nombre ni correo, solo el id, el rol y el código: ms-cursos lo usa para traducir el código institucional del CSV de matrícula a un \texttt{userId}, y ms-notificaciones para crear las preferencias por defecto. Un ítem \texttt{USERCODE\#} garantiza que el código no se repita dentro del tenant.

Todos los roles leen su tenant, los periodos y su propio perfil, y pueden cambiar solo su idioma y zona horaria (\texttt{PATCH /users/me}). Docentes y administradores pueden listar y buscar usuarios del tenant, por rol o por código, por ejemplo para armar una matrícula; solo el administrador crea e inicia periodos. Como todo se guarda en la partición de su \texttt{tenant\_id}, un usuario de otra institución recibe 404 aunque conozca el id. ms-usuarios no consume eventos.

\subsubsection{Endpoints}
\begin{table}[htbp]
\centering\footnotesize
\begin{tabularx}{\textwidth}{lXll}
\toprule
Método & Ruta & Roles & Emite \\
\midrule
GET    & \texttt{/health} & pública & --- \\
GET    & \texttt{/tenant} & todos & --- \\
GET    & \texttt{/terms} & todos & --- \\
POST   & \texttt{/terms} & ADMIN & --- \\
POST   & \texttt{/terms/\{termId\}/start} & ADMIN & \texttt{term.started} \\
POST   & \texttt{/users/me} & todos & \texttt{user.provisioned} \\
GET    & \texttt{/users/me} & todos & --- \\
PATCH  & \texttt{/users/me} & todos & --- \\
GET    & \texttt{/users} & TEACHER, ADMIN & --- \\
GET    & \texttt{/users/\{userId\}} & TEACHER, ADMIN & --- \\
\bottomrule
\end{tabularx}
\caption{Endpoints de ms-usuarios (\texttt{apis/ms-usuarios.openapi.yaml} de \texttt{nubeduca-contratos}); \textit{todos} = STUDENT, TEACHER y ADMIN.}
\end{table}

\subsubsection{Eventos}
\textbf{Publica:}
\begin{itemize}
  \item \texttt{lms.users.tenant.created.v1}
  \item \texttt{lms.users.term.started.v1}
  \item \texttt{lms.users.user.provisioned.v1}
\end{itemize}

\textbf{Consume:}
\textit{Ninguno.}


\subsubsection{Estructura JSON de la tabla \texttt{<stage>-t\_usuarios}}
Clave de partición \texttt{tenant\_id}, clave de ordenamiento \texttt{sk}.

\begin{table}[htbp]
\centering\small
\begin{tabularx}{\textwidth}{>{\raggedright\arraybackslash}p{0.42\textwidth}X}
\toprule
\texttt{sk} & Qué es \\
\midrule
\texttt{TENANT} & Institución \\
\texttt{TERM\#<termId>} & Periodo académico \\
\texttt{TERMNAME\#<name>} & Reserva del nombre del periodo en el tenant \\
\texttt{USER\#<userId>} & Perfil \\
\texttt{USERCODE\#<código>} & Código institucional $\rightarrow$ usuario (único en el tenant) \\
\texttt{ROSTER\#<email>} & Padrón: lo que la institución ya sabe de ese correo, para el primer login \\
\texttt{EVENTO\#<eventId>} & Outbox (lo lee el Stream) \\
\bottomrule
\end{tabularx}
\caption{Prefijos de \texttt{sk} en la tabla de ms-usuarios.}
\end{table}

\begin{lstlisting}
{
  "tenant_id": "utec",
  "sk": "TENANT",
  "entidad": "TENANT",
  "name": "Universidad de Ingeniería y Tecnología",
  "timeZone": "America/Lima",
  "locale": "es-PE",
  "version": 1,
  "createdAt": "2026-11-13T15:00:00Z"
}
\end{lstlisting}

\begin{lstlisting}
{
  "tenant_id": "utec",
  "sk": "TERM#trm_01JCP0KWJDXF9GJ9YNTK2H84A9",
  "entidad": "TERM",
  "termId": "trm_01JCP0KWJDXF9GJ9YNTK2H84A9",
  "name": "2026-2",
  "startsOn": "2026-08-17",
  "endsOn": "2026-12-12",
  "status": "ACTIVE",
  "version": 2,
  "createdBy": "16248d98-4e7a-41b9-a08b-4929177d2e2d",
  "createdAt": "2026-11-13T15:00:00Z",
  "startedBy": "16248d98-4e7a-41b9-a08b-4929177d2e2d",
  "startedAt": "2026-11-13T15:05:00Z"
}
\end{lstlisting}

\begin{lstlisting}
{
  "tenant_id": "utec",
  "sk": "USER#11111111-1111-4111-8111-111111111111",
  "entidad": "USER",
  "userId": "11111111-1111-4111-8111-111111111111",
  "role": "STUDENT",
  "institutionalCode": "202310457",
  "givenName": "Ana",
  "familyName": "Torres",
  "email": "alumno.demo@utec.edu.pe",
  "locale": "es-PE",
  "timeZone": "America/Lima",
  "version": 1,
  "createdAt": "2026-11-10T15:00:00Z"
}
\end{lstlisting}

\begin{lstlisting}
{
  "tenant_id": "utec",
  "sk": "ROSTER#alumno.demo@utec.edu.pe",
  "entidad": "ROSTER",
  "institutionalCode": "202310457",
  "givenName": "Ana",
  "familyName": "Torres",
  "locale": "es-PE",
  "timeZone": "America/Lima"
}
\end{lstlisting}

Formato de los ítems documentado en el README de \texttt{nubeduca-ms-usuarios}.

\subsubsection{Datos de prueba}
El script \texttt{scripts/poblar\_datos.py} crea los tres tenants (UTEC, UPC y PUCP) y publica un \texttt{tenant.created} por cada uno; si el tenant ya existe, no lo publica de nuevo. Después carga 30\,000 perfiles: UTEC 15\,000, UPC 10\,000 y PUCP 5\,000, con 90\,\% de alumnos, 9\,\% de docentes y 1\,\% de administradores, cada uno con su código institucional. Los perfiles no escriben ítems \texttt{EVENTO\#}, para no inundar el bus, y sus identificadores son fijos, así que el script es idempotente. Incluye los usuarios demo del frontend y un periodo académico activo de seed, el mismo que usa ms-cursos.

Carga en la tabla desplegada \texttt{dev-t\_usuarios} (22 de noviembre de 2026):

\begin{lstlisting}
$ python scripts/poblar_datos.py --stage dev
utec: tenant CREADO
upc: tenant CREADO
pucp: tenant CREADO
utec: 15000 usuarios
upc: 10000 usuarios
pucp: 5000 usuarios
Total: 30000 usuarios en dev-t_usuarios
\end{lstlisting}

El conteo de los perfiles (\texttt{USER\#}) por tenant en la tabla desplegada, con \texttt{aws dynamodb query ... --select COUNT}, confirma UTEC 15\,000, UPC 10\,000 y PUCP 5\,000 (figura~\ref{fig:ms-usuarios-datos}).

\evidencia{ms-usuarios-datos.png}{ms-usuarios: conteo de perfiles por tenant en \texttt{dev-t\_usuarios}.}{fig:ms-usuarios-datos}

\subsubsection{Resultados y evidencias}
\evidencia{ms-usuarios-evidencia.png}{ms-usuarios desplegado en AWS: \texttt{GET /users?role=TEACHER\&limit=3} como administrador de UTEC devuelve 200 con docentes del seed.}{fig:ms-usuarios}

\subsection{ms-cursos --- Cursos, secciones, matrícula y anuncios}
\label{sec:ms-cursos}
\textit{Responsable: Benjamín y Edinson (ACA-03, ACA-04, COL-01). Repositorio: \url{https://github.com/nubeduca/nubeduca-ms-cursos}.}

ms-cursos gestiona los cursos de cada periodo, sus secciones, quién está matriculado y los anuncios. Un curso nace en borrador (\texttt{DRAFT}), se publica (\texttt{PUBLISHED}) cuando los alumnos ya pueden verlo y se archiva (\texttt{ARCHIVED}) para dejarlo de solo lectura; cada transición emite su evento y sube la \texttt{aggregateVersion} del curso (1, 2 y 3), así los consumidores descartan los eventos que llegan en desorden. Solo se crean cursos en un periodo activo, y el código del curso (por ejemplo \texttt{CS2032}) es único dentro del periodo gracias a un ítem de reserva \texttt{COURSECODE\#} que va en la misma transacción.

El servicio no llama a ms-usuarios: mantiene dos proyecciones locales alimentadas por eventos. Con \texttt{term.started} guarda qué periodos están activos, y con \texttt{user.provisioned} guarda el rol y el código institucional de cada usuario. Con esa segunda proyección funciona la matrícula por CSV (\texttt{institutionalCode,role,section}, hasta 2\,000 filas): cada fila válida crea una matrícula y emite un \texttt{enrollment.created}, y las filas con error se devuelven con su número de línea y el motivo, sin detener la importación. Cada matrícula se guarda tres veces en una sola transacción: por curso (\texttt{ENROLL\#}, para listar los matriculados), por id (\texttt{ENROLLID\#}, para el \texttt{DELETE}) y por usuario (\texttt{USERCOURSE\#}, para los cursos del alumno), de modo que cada consulta es una \texttt{query} por prefijo. Los anuncios se ordenan por su ULID, del más reciente al más antiguo, y su evento lleva solo el título: el cuerpo se lee con \texttt{GET}.

El tenant y el rol salen siempre del token. El administrador ve y gestiona todos los cursos de su tenant. El docente crea cursos y solo puede editar, publicar, archivar, crear secciones y matricular en los suyos; para publicar un anuncio, además, tiene que estar matriculado como \texttt{TEACHER} en el curso, y el curso debe estar publicado. El alumno solo ve los cursos publicados donde está matriculado, con sus secciones y anuncios; a un curso sin acceso se le responde 404, para no revelar que existe.

\subsubsection{Endpoints}
\begin{table}[htbp]
\centering\footnotesize
\begin{tabularx}{\textwidth}{lXll}
\toprule
Método & Ruta & Roles & Emite \\
\midrule
GET    & \texttt{/health} & pública & --- \\
GET    & \texttt{/courses} & todos & --- \\
POST   & \texttt{/courses} & TEACHER, ADMIN & \texttt{course.created} \\
GET    & \texttt{/courses/\{courseId\}} & todos & --- \\
PATCH  & \texttt{/courses/\{courseId\}} & TEACHER, ADMIN & --- \\
POST   & \texttt{/courses/\{courseId\}/publish} & TEACHER, ADMIN & \texttt{course.published} \\
POST   & \texttt{/courses/\{courseId\}/archive} & TEACHER, ADMIN & \texttt{course.archived} \\
GET    & \texttt{/courses/\{courseId\}/sections} & todos & --- \\
POST   & \texttt{/courses/\{courseId\}/sections} & TEACHER, ADMIN & --- \\
GET    & \texttt{/courses/\{courseId\}/enrollments} & TEACHER, ADMIN & --- \\
POST   & \texttt{/courses/\{courseId\}/enrollments} & TEACHER, ADMIN & \texttt{enrollment.created} \\
POST   & \texttt{/courses/\{courseId\}/enrollments/import} & TEACHER, ADMIN & \texttt{enrollment.created} \\
DELETE & \texttt{/courses/\{courseId\}/enrollments/\{enrollmentId\}} & TEACHER, ADMIN & \texttt{enrollment.dropped} \\
GET    & \texttt{/courses/\{courseId\}/announcements} & todos & --- \\
POST   & \texttt{/courses/\{courseId\}/announcements} & TEACHER & \texttt{announcement.published} \\
\bottomrule
\end{tabularx}
\caption{Endpoints de ms-cursos (\texttt{apis/ms-cursos.openapi.yaml} de \texttt{nubeduca-contratos}); \textit{todos} = STUDENT, TEACHER y ADMIN.}
\end{table}

\subsubsection{Eventos}
\textbf{Publica:}
\begin{itemize}
  \item \texttt{lms.courses.course.created.v1}
  \item \texttt{lms.courses.course.published.v1}
  \item \texttt{lms.courses.course.archived.v1}
  \item \texttt{lms.courses.enrollment.created.v1}
  \item \texttt{lms.courses.enrollment.dropped.v1}
  \item \texttt{lms.courses.announcement.published.v1}
\end{itemize}

\textbf{Consume:}
\begin{itemize}
  \item \texttt{lms.users.term.started.v1}
  \item \texttt{lms.users.user.provisioned.v1}
\end{itemize}

\paragraph{Hallazgo de la prueba de carga (PERF-01).}
\texttt{GET /courses} hace \texttt{query} de toda la partición del tenant y filtra \texttt{ownerId} con
\texttt{FilterExpression}, sin índice: con el seed de 30\,000 tarda unos 1{,}7\,s (p95 1{,}84\,s)
incluso con un usuario, mientras que \texttt{GET /courses/\{courseId\}} responde en 43\,ms (p95). Un
GSI por \texttt{ownerId} lo dejaría en decenas de milisegundos (sección~\ref{sec:pruebas}).

El índice no se agregó antes de la entrega: la tabla de \texttt{serverless.yml} en \texttt{main} sigue
sin GSI, así que la latencia de \texttt{GET /courses} es la medida en PERF-01. Agregar el GSI por
\texttt{ownerId} queda como trabajo futuro (sección~\ref{sec:conclusiones}).

\subsubsection{Estructura JSON de la tabla \texttt{<stage>-t\_cursos}}
Clave de partición \texttt{tenant\_id}, clave de ordenamiento \texttt{sk}.

\begin{table}[htbp]
\centering\small
\begin{tabularx}{\textwidth}{>{\raggedright\arraybackslash}p{0.46\textwidth}X}
\toprule
\texttt{sk} & Qué es \\
\midrule
\texttt{COURSE\#<courseId>} & Curso \\
\texttt{COURSECODE\#<termId>\#<code>} & Reserva del código del curso en el periodo \\
\texttt{SECTION\#<courseId>\#<sectionId>} & Sección del curso \\
\texttt{ENROLL\#<courseId>\#<userId>} & Matrícula \\
\texttt{ENROLLID\#<enrollmentId>} & Puntero a la matrícula, para el \texttt{DELETE} \\
\texttt{USERCOURSE\#<userId>\#<courseId>} & Cursos de un usuario \\
\texttt{ANN\#<courseId>\#<announcementId>} & Anuncio \\
\texttt{PROY\#TERM\#<termId>} & Periodo activo (proyección de \texttt{term.started}) \\
\texttt{PROY\#USER\#<userId>} & Rol y código del usuario (proyección de \texttt{user.provisioned}) \\
\texttt{PROY\#USERCODE\#<código>} & Código institucional $\rightarrow$ usuario, para el CSV \\
\texttt{EVENTO\#}, \texttt{INBOX\#}, \texttt{AGG\#} & Outbox, inbox y versión de agregados de la plantilla \\
\bottomrule
\end{tabularx}
\caption{Prefijos de \texttt{sk} en la tabla de ms-cursos.}
\end{table}

\begin{lstlisting}
{
  "tenant_id": "utec",
  "sk": "COURSE#crs_00000000000000000000000020",
  "entidad": "COURSE",
  "version": 1,
  "courseId": "crs_00000000000000000000000020",
  "termId": "trm_0000000000000000000000SEED",
  "code": "SEED-00020",
  "title": "Cloud Computing 3",
  "description": "Curso de seed para pruebas de carga.",
  "status": "PUBLISHED",
  "ownerId": "3651db88-79cd-5df8-bd8c-488941b6e878",
  "createdAt": "2026-08-17T13:00:00Z",
  "seed": true
}
\end{lstlisting}
Ítem real leído de la tabla desplegada \texttt{dev-t\_cursos} con \texttt{aws dynamodb get-item} (22-Nov-2026).

Los demás tipos de ítem muestran los atributos que escribe el código de \texttt{api/} y
\texttt{consumidores/}, con un curso, un docente y la alumna demo Ana Torres de ejemplo. La reserva
del código y el puntero de la matrícula:

\begin{lstlisting}
{
  "tenant_id": "utec",
  "sk": "COURSECODE#trm_01JCP0KWJDXF9GJ9YNTK2H84A9#CS2032",
  "entidad": "COURSECODE",
  "courseId": "crs_01JCQRZHQT9F4WY2JYPDV546K7"
}
\end{lstlisting}

\begin{lstlisting}
{
  "tenant_id": "utec",
  "sk": "ENROLLID#enr_01JCR1B7Y2K8M4Q6T9V3W5X0Z1",
  "entidad": "ENROLLMENT_ID",
  "courseId": "crs_01JCQRZHQT9F4WY2JYPDV546K7",
  "userId": "11111111-1111-4111-8111-111111111111"
}
\end{lstlisting}

Sección, matrícula y su copia por usuario:

\begin{lstlisting}
{
  "tenant_id": "utec",
  "sk": "SECTION#crs_01JCQRZHQT9F4WY2JYPDV546K7#sec_01JCR0A4T7Q9W2N8M5K3H6G1F0",
  "entidad": "SECTION",
  "courseId": "crs_01JCQRZHQT9F4WY2JYPDV546K7",
  "sectionId": "sec_01JCR0A4T7Q9W2N8M5K3H6G1F0",
  "name": "Sección 1",
  "createdBy": "c62d9930-b473-4b1d-adec-72fdda7cf599",
  "createdAt": "2026-11-13T15:10:00Z"
}
\end{lstlisting}

\begin{lstlisting}
{
  "tenant_id": "utec",
  "sk": "ENROLL#crs_01JCQRZHQT9F4WY2JYPDV546K7#11111111-1111-4111-8111-111111111111",
  "entidad": "ENROLLMENT",
  "courseId": "crs_01JCQRZHQT9F4WY2JYPDV546K7",
  "version": 1,
  "enrollmentId": "enr_01JCR1B7Y2K8M4Q6T9V3W5X0Z1",
  "userId": "11111111-1111-4111-8111-111111111111",
  "role": "STUDENT",
  "sectionId": "sec_01JCR0A4T7Q9W2N8M5K3H6G1F0",
  "createdAt": "2026-11-14T16:00:00Z"
}
\end{lstlisting}

\begin{lstlisting}
{
  "tenant_id": "utec",
  "sk": "USERCOURSE#11111111-1111-4111-8111-111111111111#crs_01JCQRZHQT9F4WY2JYPDV546K7",
  "entidad": "USER_COURSE",
  "courseId": "crs_01JCQRZHQT9F4WY2JYPDV546K7",
  "userId": "11111111-1111-4111-8111-111111111111",
  "role": "STUDENT",
  "enrollmentId": "enr_01JCR1B7Y2K8M4Q6T9V3W5X0Z1"
}
\end{lstlisting}

Anuncio:

\begin{lstlisting}
{
  "tenant_id": "utec",
  "sk": "ANN#crs_01JCQRZHQT9F4WY2JYPDV546K7#ann_01JCS2D9F4H6K8N0P2R4T6V8X0",
  "entidad": "ANNOUNCEMENT",
  "courseId": "crs_01JCQRZHQT9F4WY2JYPDV546K7",
  "announcementId": "ann_01JCS2D9F4H6K8N0P2R4T6V8X0",
  "title": "Bienvenidos al curso",
  "body": "El primer laboratorio es el lunes.",
  "authorId": "c62d9930-b473-4b1d-adec-72fdda7cf599",
  "publishedAt": "2026-11-15T14:00:00Z"
}
\end{lstlisting}

Proyecciones de los eventos de ms-usuarios:

\begin{lstlisting}
{
  "tenant_id": "utec",
  "sk": "PROY#TERM#trm_01JCP0KWJDXF9GJ9YNTK2H84A9",
  "entidad": "PROYECCION",
  "termId": "trm_01JCP0KWJDXF9GJ9YNTK2H84A9",
  "name": "2026-2",
  "startsOn": "2026-08-17",
  "endsOn": "2026-12-12",
  "status": "ACTIVE",
  "recibidoDe": "lms.users",
  "version": 2
}
\end{lstlisting}

\begin{lstlisting}
{
  "tenant_id": "utec",
  "sk": "PROY#USER#11111111-1111-4111-8111-111111111111",
  "entidad": "PROYECCION",
  "userId": "11111111-1111-4111-8111-111111111111",
  "role": "STUDENT",
  "institutionalCode": "202310457",
  "recibidoDe": "lms.users",
  "version": 1
}
\end{lstlisting}

\begin{lstlisting}
{
  "tenant_id": "utec",
  "sk": "PROY#USERCODE#202310457",
  "entidad": "PROYECCION",
  "userId": "11111111-1111-4111-8111-111111111111",
  "role": "STUDENT"
}
\end{lstlisting}

\subsubsection{Datos de prueba}
El 20-Nov-2026, en la cuenta de Dev 2 (stage \texttt{dev}), \texttt{poblar\_datos.py}
cargó \textbf{30\,000 cursos} en \texttt{dev-t\_cursos}: UTEC 15\,000, UPC 10\,000 y PUCP 5\,000. Cada tenant
recibe un periodo activo de seed y 20 docentes de seed; los estados quedan en 60\,\% PUBLISHED, 30\,\% DRAFT y
10\,\% ARCHIVED. Un conteo de los ítems \texttt{COURSE\#} de UTEC devolvió 15\,000.

El 22-Nov-2026 se repitió la carga en otra cuenta (stage \texttt{dev}, cuenta de Dev 4):

\begin{lstlisting}
$ python scripts/poblar_datos.py --stage dev
utec: 15000 cursos
upc: 10000 cursos
pucp: 5000 cursos
Total: 30000 cursos en dev-t_cursos
\end{lstlisting}

El conteo de los ítems \texttt{COURSE\#} de los tres tenants, con \texttt{aws dynamodb query ... --select COUNT}, confirma UTEC 15\,000, UPC 10\,000 y PUCP 5\,000 (figura~\ref{fig:ms-cursos-datos}).

\evidencia{ms-cursos-datos.png}{ms-cursos: conteo de cursos por tenant en \texttt{dev-t\_cursos}.}{fig:ms-cursos-datos}

\subsubsection{Resultados y evidencias}
\evidencia{ms-cursos-evidencia.png}{ms-cursos desplegado en AWS: \texttt{GET /courses/\{courseId\}} como administrador de UTEC devuelve 200 con el primer curso del seed.}{fig:ms-cursos}

\subsection{ms-contenidos --- Archivos (S3), módulos e ítems}
\label{sec:ms-contenidos}
\textit{Responsable: Alexander (CON-01, CON-02). Repositorio: \url{https://github.com/nubeduca/nubeduca-ms-contenidos}.}

\pendiente{Dos o tres párrafos: qué resuelve este microservicio, decisiones de diseño y qué roles (STUDENT, TEACHER, ADMIN) usan cada operación.}

\subsubsection{Endpoints}
\pendiente{Tabla con método, ruta, roles y evento que emite. Sale del OpenAPI del microservicio en la carpeta \texttt{apis/} de \texttt{nubeduca-contratos} (campos \texttt{x-roles} y \texttt{x-emits}).}

\subsubsection{Eventos}
\textbf{Publica:}
\begin{itemize}
  \item \texttt{lms.content.file.scanned.v1}
  \item \texttt{lms.content.module.published.v1}
\end{itemize}

\textbf{Consume:}
\begin{itemize}
  \item \texttt{lms.courses.course.created.v1}
  \item \texttt{lms.courses.enrollment.created.v1}
  \item \texttt{lms.courses.enrollment.dropped.v1}
\end{itemize}

\pendiente{Alexander: explicar el flujo de archivos (URL prefirmada, evento de S3, Lambda de antivirus, estados CLEAN/INFECTED) y el bucket usado.}

\subsubsection{Estructura JSON de la tabla \texttt{<stage>-t\_contenidos}}
Clave de partición \texttt{tenant\_id}, clave de ordenamiento \texttt{sk}.

\begin{lstlisting}
{
  "tenant_id": "utec",
  "sk": "<PREFIJO>#<id>",
  "...": "completar con un ítem real de la tabla"
}
\end{lstlisting}
\pendiente{Reemplazar por un ítem real de cada tipo que guarda la tabla (por ejemplo con \texttt{aws dynamodb scan --table-name dev-t\_contenidos --max-items 1}). Una tabla con los prefijos de \texttt{sk} y un bloque JSON por tipo de ítem.}

\subsubsection{Datos de prueba}
\pendiente{Evidencia de los 30\,000 registros (UTEC 15\,000, UPC 10\,000, PUCP 5\,000): salida de \texttt{poblar\_datos.py} o un conteo por tenant con \texttt{aws dynamodb query ... --select COUNT}.}

\subsubsection{Resultados y evidencias}
\evidencia{ms-contenidos-evidencia.png}{ms-contenidos: prueba funcional (por ejemplo, una llamada en Postman o una captura del frontend).}{fig:ms-contenidos}

\subsection{ms-evaluaciones --- Tareas y entregas}
\label{sec:ms-evaluaciones}
\textit{Responsable: Edinson (EVA-01). Repositorio: \url{https://github.com/nubeduca/nubeduca-ms-evaluaciones}.}

\pendiente{Dos o tres párrafos: qué resuelve este microservicio, decisiones de diseño y qué roles (STUDENT, TEACHER, ADMIN) usan cada operación.}

\subsubsection{Endpoints}
\begin{table}[htbp]
\centering\footnotesize
\begin{tabularx}{\textwidth}{lXll}
\toprule
Método & Ruta & Roles & Emite \\
\midrule
GET    & \texttt{/health} & pública & --- \\
GET    & \texttt{/courses/\{courseId\}/assignments} & todos & --- \\
POST   & \texttt{/courses/\{courseId\}/assignments} & TEACHER & --- \\
GET    & \texttt{/assignments/\{assignmentId\}} & todos & --- \\
PATCH  & \texttt{/assignments/\{assignmentId\}} & TEACHER & --- \\
POST   & \texttt{/assignments/\{assignmentId\}/publish} & TEACHER & \texttt{assignment.published} \\
GET    & \texttt{/assignments/\{assignmentId\}/submissions} & STUDENT, TEACHER & --- \\
POST   & \texttt{/assignments/\{assignmentId\}/submissions} & STUDENT & \texttt{submission.submitted} \\
GET    & \texttt{/submissions/\{submissionId\}} & STUDENT, TEACHER & --- \\
POST   & \texttt{/submissions/\{submissionId\}/grade} & TEACHER & \texttt{submission.graded} \\
\bottomrule
\end{tabularx}
\caption{Endpoints de ms-evaluaciones (\texttt{apis/ms-evaluaciones.openapi.yaml} de \texttt{nubeduca-contratos}); \textit{todos} = STUDENT, TEACHER y ADMIN.}
\end{table}

\subsubsection{Eventos}
\textbf{Publica:}
\begin{itemize}
  \item \texttt{lms.assessment.assignment.published.v1}
  \item \texttt{lms.assessment.submission.submitted.v1}
  \item \texttt{lms.assessment.submission.graded.v1}
\end{itemize}

\textbf{Consume:}
\begin{itemize}
  \item \texttt{lms.courses.enrollment.created.v1}
  \item \texttt{lms.courses.enrollment.dropped.v1}
  \item \texttt{lms.content.file.scanned.v1}
\end{itemize}


\subsubsection{Estructura JSON de la tabla \texttt{<stage>-t\_evaluaciones}}
Clave de partición \texttt{tenant\_id}, clave de ordenamiento \texttt{sk}.

\begin{table}[htbp]
\centering\small
\begin{tabularx}{\textwidth}{>{\raggedright\arraybackslash}p{0.42\textwidth}X}
\toprule
\texttt{sk} & Qué es \\
\midrule
\texttt{ASSIGN\#<assignmentId>} & Tarea \\
\texttt{COURSEASSIGN\#<courseId>\#<dueAt>\#<assignmentId>} & Índice: tareas de un curso por fecha de entrega \\
\texttt{SUBMIT\#<submissionId>} & Entrega \\
\texttt{ASSIGNSUBMIT\#<assignmentId>\#<studentId>\#<submissionId>} & Índice: entregas de una tarea (y de un alumno en esa tarea) \\
\texttt{ATTEMPTS\#<assignmentId>\#<studentId>} & Contador de entregas del alumno en la tarea (\texttt{total}) \\
\texttt{PROY\#ENROLL\#<courseId>\#<userId>} & Matrícula, proyección de \texttt{enrollment.created} \\
\texttt{PROY\#FILE\#<fileId>} & Archivo escaneado, proyección de \texttt{file.scanned} \\
\texttt{EVENTO\#}, \texttt{INBOX\#}, \texttt{AGG\#} & Outbox, inbox y versión de agregados de la plantilla \\
\bottomrule
\end{tabularx}
\caption{Prefijos de \texttt{sk} en la tabla de ms-evaluaciones.}
\end{table}

\begin{lstlisting}
{
  "tenant_id": "utec",
  "sk": "ASSIGN#asg_01JCF4SQ7XAMY8HTMRFHAWZ73F",
  "entidad": "ASSIGNMENT",
  "assignmentId": "asg_01JCF4SQ7XAMY8HTMRFHAWZ73F",
  "courseId": "crs_01JCQRZHQT9F4WY2JYPDV546K7",
  "title": "Laboratorio 4: SQS y DLQ",
  "instructions": "Sube tu informe en PDF.",
  "dueAt": "2026-11-20T23:59:00Z",
  "maxScore": 20,
  "allowLate": false,
  "status": "PUBLISHED",
  "ownerId": "c62d9930-b473-4b1d-adec-72fdda7cf599",
  "createdAt": "2026-11-11T15:00:00Z",
  "version": 2
}
\end{lstlisting}

\begin{lstlisting}
{
  "tenant_id": "utec",
  "sk": "SUBMIT#sbm_01JCY7SYF5X0ZHRJ7KQSVX4BP0",
  "entidad": "SUBMISSION",
  "submissionId": "sbm_01JCY7SYF5X0ZHRJ7KQSVX4BP0",
  "assignmentId": "asg_01JCF4SQ7XAMY8HTMRFHAWZ73F",
  "courseId": "crs_01JCQRZHQT9F4WY2JYPDV546K7",
  "studentId": "a18577e2-e005-4a49-a29e-755dbf08f5b6",
  "attempt": 1,
  "fileIds": ["fil_01JC8P2K48BKYRT5ATH3DZDA66"],
  "comment": "Adjunto mi informe.",
  "late": false,
  "status": "GRADED",
  "submittedAt": "2026-11-12T15:22:00Z",
  "score": 17,
  "feedback": "Buen trabajo.",
  "graderId": "c62d9930-b473-4b1d-adec-72fdda7cf599",
  "gradedAt": "2026-11-13T15:23:00Z",
  "version": 2
}
\end{lstlisting}

\begin{lstlisting}
{
  "tenant_id": "utec",
  "sk": "PROY#ENROLL#crs_01JCQRZHQT9F4WY2JYPDV546K7#a18577e2-e005-4a49-a29e-755dbf08f5b6",
  "entidad": "PROYECCION",
  "enrollmentId": "enr_01JC5F0KPBVF7HK1VV8H23JGMW",
  "courseId": "crs_01JCQRZHQT9F4WY2JYPDV546K7",
  "userId": "a18577e2-e005-4a49-a29e-755dbf08f5b6",
  "role": "STUDENT",
  "recibidoDe": "lms.courses",
  "version": 1
}
\end{lstlisting}

\begin{lstlisting}
{
  "tenant_id": "utec",
  "sk": "PROY#FILE#fil_01JC8P2K48BKYRT5ATH3DZDA66",
  "entidad": "PROYECCION",
  "fileId": "fil_01JC8P2K48BKYRT5ATH3DZDA66",
  "result": "CLEAN",
  "purpose": "SUBMISSION",
  "ownerId": "a18577e2-e005-4a49-a29e-755dbf08f5b6",
  "recibidoDe": "lms.content",
  "version": 1
}
\end{lstlisting}

Formato de los ítems documentado en el README de \texttt{nubeduca-ms-evaluaciones}.

\subsubsection{Datos de prueba}
\pendiente{Evidencia de los 30\,000 registros (UTEC 15\,000, UPC 10\,000, PUCP 5\,000): salida de \texttt{poblar\_datos.py} o un conteo por tenant con \texttt{aws dynamodb query ... --select COUNT}.}

\subsubsection{Resultados y evidencias}
\evidencia{ms-evaluaciones-evidencia.png}{ms-evaluaciones: prueba funcional (por ejemplo, una llamada en Postman o una captura del frontend).}{fig:ms-evaluaciones}

\subsection{ms-calificaciones --- Libro de notas y auditoría}
\label{sec:ms-calificaciones}
\textit{Responsable: Edinson y Benjamín (EVA-02); Alexander (DAT-01). Repositorio: \url{https://github.com/nubeduca/nubeduca-ms-calificaciones}.}

\pendiente{Dos o tres párrafos: qué resuelve este microservicio, decisiones de diseño y qué roles (STUDENT, TEACHER, ADMIN) usan cada operación.}

\subsubsection{Endpoints}
\pendiente{Tabla con método, ruta, roles y evento que emite. Sale del OpenAPI del microservicio en la carpeta \texttt{apis/} de \texttt{nubeduca-contratos} (campos \texttt{x-roles} y \texttt{x-emits}).}

\subsubsection{Eventos}
\textbf{Publica:}
\begin{itemize}
  \item \texttt{lms.grading.grade.changed.v1}
  \item \texttt{lms.grading.grade.posted.v1 (bus + SNS FIFO)}
\end{itemize}

\textbf{Consume:}
\begin{itemize}
  \item \texttt{lms.assessment.assignment.published.v1}
  \item \texttt{lms.courses.enrollment.created.v1}
  \item \texttt{lms.courses.enrollment.dropped.v1}
  \item \texttt{lms.assessment.submission.graded.v1}
\end{itemize}

\pendiente{Explicar la publicación de notas en el tópico SNS FIFO \texttt{nubeduca-<stage>-notas.fifo} (\texttt{MessageGroupId = <tenantId>:<courseId>}) y la auditoría de \texttt{grade.changed} (DAT-01).}

\subsubsection{Estructura JSON de la tabla \texttt{<stage>-t\_calificaciones}}
Clave de partición \texttt{tenant\_id}, clave de ordenamiento \texttt{sk}.

\begin{lstlisting}
{
  "tenant_id": "utec",
  "sk": "<PREFIJO>#<id>",
  "...": "completar con un ítem real de la tabla"
}
\end{lstlisting}
\pendiente{Reemplazar por un ítem real de cada tipo que guarda la tabla (por ejemplo con \texttt{aws dynamodb scan --table-name dev-t\_calificaciones --max-items 1}). Una tabla con los prefijos de \texttt{sk} y un bloque JSON por tipo de ítem.}

\subsubsection{Datos de prueba}
\pendiente{Evidencia de los 30\,000 registros (UTEC 15\,000, UPC 10\,000, PUCP 5\,000): salida de \texttt{poblar\_datos.py} o un conteo por tenant con \texttt{aws dynamodb query ... --select COUNT}.}

\subsubsection{Resultados y evidencias}
\evidencia{ms-calificaciones-evidencia.png}{ms-calificaciones: prueba funcional (por ejemplo, una llamada en Postman o una captura del frontend).}{fig:ms-calificaciones}

\subsection{ms-notificaciones --- Avisos por correo}
\label{sec:ms-notificaciones}
\textit{Responsable: Guillermo (ENG-01). Repositorio: \url{https://github.com/nubeduca/nubeduca-ms-notificaciones}.}

\pendiente{Dos o tres párrafos: qué resuelve este microservicio, decisiones de diseño y qué roles (STUDENT, TEACHER, ADMIN) usan cada operación.}

\subsubsection{Endpoints}
\pendiente{Tabla con método, ruta, roles y evento que emite. Sale del OpenAPI del microservicio en la carpeta \texttt{apis/} de \texttt{nubeduca-contratos} (campos \texttt{x-roles} y \texttt{x-emits}).}

\subsubsection{Eventos}
\textbf{Publica:}
\textit{Ninguno.}

\textbf{Consume:}
\begin{itemize}
  \item \texttt{lms.users.user.provisioned.v1}
  \item \texttt{lms.courses.course.published.v1}
  \item \texttt{lms.courses.enrollment.*}
  \item \texttt{lms.courses.announcement.published.v1}
  \item \texttt{lms.content.module.published.v1}
  \item \texttt{lms.assessment.assignment.published.v1}
  \item \texttt{lms.assessment.submission.submitted.v1}
  \item \texttt{lms.grading.grade.posted.v1 (desde su cola SQS FIFO, no por regla del bus)}
\end{itemize}

\pendiente{Guillermo: explicar las preferencias por usuario y el envío de correo con SNS (confirmación de suscripción o correos verificados).}

\subsubsection{Estructura JSON de la tabla \texttt{<stage>-t\_notificaciones}}
Clave de partición \texttt{tenant\_id}, clave de ordenamiento \texttt{sk}.

\begin{lstlisting}
{
  "tenant_id": "utec",
  "sk": "<PREFIJO>#<id>",
  "...": "completar con un ítem real de la tabla"
}
\end{lstlisting}
\pendiente{Reemplazar por un ítem real de cada tipo que guarda la tabla (por ejemplo con \texttt{aws dynamodb scan --table-name dev-t\_notificaciones --max-items 1}). Una tabla con los prefijos de \texttt{sk} y un bloque JSON por tipo de ítem.}

\subsubsection{Datos de prueba}
\pendiente{Evidencia de los 30\,000 registros (UTEC 15\,000, UPC 10\,000, PUCP 5\,000): salida de \texttt{poblar\_datos.py} o un conteo por tenant con \texttt{aws dynamodb query ... --select COUNT}.}

\subsubsection{Resultados y evidencias}
\evidencia{ms-notificaciones-evidencia.png}{ms-notificaciones: prueba funcional (por ejemplo, una llamada en Postman o una captura del frontend).}{fig:ms-notificaciones}

